\subsection{Coxeter 变换的定义}

相对于 W 的有序室是指由 H 所确定的一个室 C 与一个双射 \(i \mapsto H_{i}\)（从 \(\{1, 2, \ldots, l\}\) 到 C 的墙集合上）所组成的偶（cf.~§3, no. 9, Prop. 7）。

定义 1. 由有序室 \((\mathrm{C},(\mathrm{H}_i)_{1\leqslant i\leqslant l})\) 确定的 Coxeter 变换是 W 的元素 \(c = s_{\mathrm{H}_1}s_{\mathrm{H}_2}\dots s_{\mathrm{H}_l}\)。

命题 1. 所有 Coxeter 变换在 W 中彼此共轭。

由于 W 传递地置换由 H 确定的诸室（§3, no. 2, Th. 1），我们归结为证明下述结论：设 \((\mathrm{C},(\mathrm{H}_{i})_{1\leqslant i\leqslant l})\) 是一个有序室，\(\pi\in S_{l}\)；则 \(s_{H_{1}}s_{H_{2}}\ldots s_{H_{l}}\) 与 \(s_{H_{\pi(1)}}s_{H_{\pi(2)}}\ldots s_{H_{\pi(l)}}\) 在 W 中共轭。考虑到 §4, no. 8, Prop. 8，这将由下面的引理立即得出：

引理 1. 设 X 是有限森林，\(x \mapsto g_{x}\) 是从 X 到群 \(\Gamma\) 的一个映射，使得只要 x 与 y 在 X 中不相连，\(g_{x}\) 与 \(g_{y}\) 就共轭。设 T 是 X 上全体全序所成的集合。对每个 \(\xi \in T\)，用 \(p_{\xi}\) 表示 \(\Gamma\) 中将序列 \((g_{x})_{x \in X}\) 按 \(\xi\) 排列所得到的乘积。则诸元素 \(p_{\xi}\) 在 \(\Gamma\) 中共轭。

\begin{enumerate}
\def\labelenumi{\arabic{enumi})}
\item
  我们对 \(n = \operatorname{Card} X\) 作归纳。\(n = 1\) 的情形是直接的，故设 \(n \geqslant 2\)。\(X\) 中存在一个端点 \(a\)（Chap. IV, Appendix, no. 3, Prop. 2）。设 \(b \in X - \{a\}\) 是与 \(a\) 相连的顶点（若这样的顶点存在）；若 \(a\) 不与 \(X - \{a\}\) 中的任何顶点相连，则 \(b\) 取 \(X - \{a\}\) 中任意顶点。在所有情形下，\(g_a\) 与 \(g_x\) 对 \(x \neq b\) 都交换。取 \(\eta \in T\) 使 \(a\) 为 \(X\) 的最大元且 \(b\) 为 \(X - \{a\}\) 的最大元；任取 \(\xi \in T\)，证明 \(p_\xi, p_\eta\) 共轭。
\item
  首先假设对 \(\xi\) 而言，\(a\) 是 X 的最大元且 \(b\) 是 \(\mathrm{X} - \{a\}\) 的最大元。设 \(\mathrm{X}'\) 是全子图 \(\mathrm{X} - \{a\}\)，它是一个森林。定义映射 \(x \mapsto g_x'\)（由 \(\mathrm{X}'\) 到 \(\Gamma\)）如下：置 \(g_x' = g_x\)（若 \(x \neq b\)），而 \(g_b' = g_b g_a\)。设 \(\xi', \eta'\) 是 \(\xi, \eta\) 在 \(\mathrm{X}'\) 上的限制。归纳假设可以应用，故 \(p_{\xi'}\) 与 \(p_{\eta'}\) 共轭。但显然 \(p_{\xi'} = p_{\xi}\)，\(p_{\eta'} = p_{\eta}\)，引理在此情形得证。
\item
  假设 \(a\) 是 \(X\) 对 \(\xi\) 而言的最大元。设 \(X_1\)（相应地，\(X_2\)）是 \(X - \{a\}\) 中严格大于（相应地，小于）\(b\) 的元素所成的集合；设 \(\xi_i\) 是 \(\xi\) 在 \(X_i\) 上的限制。则
\end{enumerate}

\[
p _ {\xi} = p _ {\xi_ {1}} g _ {b} p _ {\xi_ {2}} g _ {a} = p _ {\xi_ {1}} g _ {b} g _ {a} p _ {\xi_ {2}},
\]

而这个元素与 \(p_{\xi_2}p_{\xi_1}g_b g_a\) 共轭。于是归结为情形 2)。

\begin{enumerate}
\def\labelenumi{\arabic{enumi})}
\setcounter{enumi}{3}
\tightlist
\item
  在一般情形，设 \(\mathrm{X}_3\)（相应地，\(\mathrm{X}_4\)）是 \(\mathbf{X}\) 中严格大于（相应地，小于）\(a\) 的元素所成的集合；设 \(\xi_i\) 是 \(\xi\) 在 \(\mathrm{X}_i\) 上的限制。则 \(p_{\xi} = p_{\xi_3}g_a p_{\xi_4}\)，而这个元素与 \(p_{\xi_4}p_{\xi_3}g_a\) 共轭。于是归结为情形 3)。
\end{enumerate}

由 Prop. 1 可知，所有 Coxeter 变换都具有同一个阶 \(h = h(\mathrm{W})\)。这个数称为 W 的 Coxeter 数。

注。设 \(W_{1},\ldots,W_{m}\) 是作用于诸空间 \(V_{1},\ldots,V_{m}\) 中的本质有限群，且都由反射生成。设 \(C_{j}\) 是相对于 \(W_{j}\) 的一个室。设 W 是群 \(W_{1}\times\cdots\times W_{m}\)，它作用于空间 \(V_{1}\times\cdots\times V_{m}\)。则 \(C_{1}\times\cdots\times C_{m}\) 是相对于 W 的一个室。由 C 确定的 W 的 Coxeter 变换是形如 \(c_{1}c_{2}\ldots c_{m}\) 的乘积，其中 \(c_{j}\) 是 \(W_{j}\) 的一个由 \(C_{j}\) 确定的 Coxeter 变换。

